\documentclass[12pt]{beamer}

\usepackage[language=british,
            sortlocale=en-GB,
            backend=biber,
            dashed=false,
            citestyle=authoryear-comp,
            bibstyle=authoryear,
            doi=true,
            useprefix=true,
            sorting=ynt,
            maxbibnames=9,
            maxcitenames=2]{biblatex}

\usepackage{filecontents}

\begin{filecontents*}{refs.bib}

@phdthesis{Bartos1999,
	Author = {Bartos, Huba},
	Location = {Budapest},
	Title = {{Morfoszintaxis {\'e}s interpret{\'a}ci{\'o}: A magyar inflexi{\'o}s jelens{\'e}gek szintaktikai h{\'a}ttere [Morphosyntax and interpretation: the syntactic background of Hungarian inflection]}},
	Year = {1999}}

@article{Coppock2013P,
	Author = {Coppock, Elizabeth},
	Journaltitle = {Natural Language Semantics},
	Pages = {345-371},
	Title = {{A semantic solution to the problem of Hungarian object agreement}},
	Volume = {21},
	Year = {2013}}

@article{Dekany2015,
	Author = {D{\'e}k{\'a}ny, {\'E}va},
	Doi = {10.1007/s11049-014-9278-0},
	Journaltitle = {Natural Language \& Linguistic Theory},
	Title = {The syntax of anaphoric possessives in Hungarian},
	Year = {2015}}

@incollection{denDikken1999,
	Author = {den Dikken, Marcel},
	Booktitle = {Crossing Boundaries: {Advances} in the Theory of Central and {Eastern European} Languages},
	Editor = {Kenesei, Istv{\'a}n},
	Location = {Amsterdam/Philadelphia},
	Pages = {137-178},
	Publisher = {John Benjamins},
	Title = {On the Structural Representation of Possession and Agreement},
	Year = {1999}}

@collection{HoeKli2008,
	Editor = {M{\"u}ller, Henrik H{\o}eg and Klinge, Alex},
	Location = {Amsterdam/Philadelphia},
	Publisher = {John Benjamins},
	Subtitle = {From morphology to discourse management},
	Title = {Essays on Nominal Determination},
	Year = {2008}}

@incollection{Szabolcsi1994,
	Author = {Szabolcsi, Anna},
	Crossref = {KieKis1994},
	Pages = {179-274},
	Title = {The Noun Phrase}}

@collection{KieKis1994,
	Editor = {Kiefer, Ferenc and {\'E}. Kiss, Katalin},
	Location = {New York},
	Publisher = {Academic Press},
	Series = {Syntax and Semantics},
	Title = {The Syntactic Structure of {Hungarian}},
	Volume = {27},
	Year = {1994}}

@article{Bernstein2008,
	Author = {Bernstein, Judy B.},
	Doi = {10.1111/j.1749-818x.2008.00091.x},
	Journaltitle = {Language and Linguistics Compass},
	Number = {6},
	Pages = {1246-1270},
	Title = {Reformulating the Determiner Phrase Analysis},
	Volume = {2},
	Year = {2008}}

@incollection{EKiss2010b,
	Author = {{\'E}. Kiss, Katalin},
	Booktitle = {Information Structure},
	Crossref = {ZimFer2010},
	Location = {Oxford/New York},
	Pages = {64-88},
	Publisher = {Oxford University Press},
	Subtitle = {Theoretical, Typological and Experimental Perspectives},
	Title = {Structural focus and exhaustivity},
	Year = {2010}}

@incollection{Longobardi2008,
	Author = {Longobardi, Giuseppe},
	Crossref = {HoeKli2008},
	Location = {Amster\-dam/Phila\-delphia},
	Pages = {189-211},
	Publisher = {John Benjamins},
	Title = {Reference to individuals, person, and the variety of mapping parameters}}

\end{filecontents*}

\addbibresource{refs.bib}

%\addbibresource{legenda.bib}

\renewcommand\nameyeardelim{\space}
\renewcommand\postnotedelim{\addcolon\space}
\DeclareFieldFormat{pages}{#1}
\DeclareFieldFormat{postnote}{#1} % no postnote prefix in "normal" citation commands
\AtEveryBibitem{%
    \clearfield{labelmonth}%
    \clearfield{issn}%
}

\usepackage{xpatch}

\xpatchbibmacro{name:andothers}{%
  \bibstring{andothers}%
}{%
  \bibstring[\emph]{andothers}%
}{}{}

% No In: for journals

\renewbibmacro{in:}{\ifentrytype{article}{}{\printtext{\bibstring{in}\intitlepunct}}}

% Lower case for prefix in bibliography
% \AtBeginBibliograph{...} is important, otherwise it breaks \Textcite{...}

\AtBeginBibliography{%
\renewbibmacro*{name:last-first}[4]{%
  \ifuseprefix%
    {\usebibmacro{name:delim}{#3#1}%
     \usebibmacro{name:hook}{#3#1}%
     \ifblank{#3}{}{%
       %\ifcapital% comment out
       %  {\mkbibnameprefix{\MakeCapital{#3}}\isdot}% and get rid of this
         {\mkbibnameprefix{#3}\isdot}%
       \ifpunctmark{'}{}{\bibnamedelimc}}%
     \mkbibnamelast{#1}\isdot%
     \ifblank{#4}{}{\bibnamedelimd\mkbibnameaffix{#4}\isdot}%
     \ifblank{#2}{}{\revsdnamepunct\bibnamedelimd\mkbibnamefirst{#2}\isdot}}
    {\usebibmacro{name:delim}{#1}%
     \usebibmacro{name:hook}{#1}%
     \mkbibnamelast{#1}\isdot%
     \ifblank{#4}{}{\bibnamedelimd\mkbibnameaffix{#4}\isdot}%
     \ifblank{#2#3}{}{\revsdnamepunct}%
     \ifblank{#2}{}{\bibnamedelimd\mkbibnamefirst{#2}\isdot}%
     \ifblank{#3}{}{\bibnamedelimd\mkbibnameprefix{#3}\isdot}}}

\renewbibmacro*{name:last}[4]{%
  \ifuseprefix%
    {\usebibmacro{name:delim}{#3#1}%
     \usebibmacro{name:hook}{#3#1}%
     \ifblank{#3}
       {}
       {%\ifcapital%<-- commented this out
        %  {\mkbibnameprefix{\MakeCapital{#3}}\isdot}% and this
          {\mkbibnameprefix{#3}\isdot}%
        \ifpunctmark{'}{}{\bibnamedelimc}}}
    {\usebibmacro{name:delim}{#1}%
     \usebibmacro{name:hook}{#1}}%
  \mkbibnamelast{#1}\isdot}%
}

% : as subtitlepunct

\renewcommand{\subtitlepunct}{\addcolon\addspace}

% Sentence case in certain fields

\DeclareFieldFormat{sentencecase}{\MakeSentenceCase{#1}}

\renewbibmacro*{title}{%
\ifthenelse{\iffieldundef{title}\AND\iffieldundef{subtitle}}
{}
{\ifthenelse{\ifentrytype{article}\OR\ifentrytype{inbook}%
  \OR\ifentrytype{incollection}\OR\ifentrytype{inproceedings}%
  \OR\ifentrytype{phdthesis}\OR\ifentrytype{book}}
  {\printtext[title]{%
    \printfield[sentencecase]{title}%
    \setunit{\subtitlepunct}%
    \printfield[sentencecase]{subtitle}}}%
  {\printtext[title]{%
    \printfield[titlecase]{title}%
    \setunit{\subtitlepunct}%
    \printfield[titlecase]{subtitle}}}%
 \newunit}%
\printfield{titleaddon}}

%%% no navigation icons

\setbeamertemplate{navigation symbols}{}

%%% no bibliography icons?

\setbeamertemplate{bibliography item}{}

%% stepwise uncovering

%\beamerdefaultoverlayspecification{<+->}

%%% other colors

\definecolor{camcosdark}{RGB}{0,145,193}
\definecolor{camcosblue}{RGB}{0,163,218}
\setbeamercolor*{section in head/foot}{parent=palette quaternary,fg=white,bg=camcosblue}%!90!black}
\setbeamercolor*{subsection in head/foot}{parent=palette secondary,fg=white,bg=camcosblue}
\setbeamercolor{structure}{bg=white,fg=black}
\setbeamercolor*{palette primary}{use=structure,bg=white,fg=black}
\setbeamercolor*{palette secondary}{use=structure,fg=white,bg=camcosblue}%!90!black}
\setbeamercolor*{palette tertiary}{use=structure,fg=white,bg=camcosblue}%!90!black}
\setbeamercolor*{palette quaternary}{use=structure,fg=white,bg=camcosblue}%!90!black}

\setbeamercolor*{alerted text}{fg=camcosblue}%!90!black}
\setbeamercolor*{part in toc}{fg=camcosblue}
\setbeamercolor*{section in toc}{fg=camcosblue}
\setbeamercolor*{item}{fg=camcosblue}

\setbeamertemplate{headline}
{%
\leavevmode
\hbox{\begin{beamercolorbox}[wd=1.00\paperwidth, ht=3.0ex, dp=1.5ex, right, rightskip=.5cm]{section in head/foot}%
        \usebeamerfont{title in head/foot}\scriptsize{\insertsection{}}
    \end{beamercolorbox}

    \begin{beamercolorbox}[wd=.00\paperwidth, ht=2.5ex, dp=1ex, left, leftskip=.5cm]{subsection in head/foot}%
        \usebeamerfont{title in head/foot}%\insertsubsection
    \end{beamercolorbox}
    }
}

\setbeamertemplate{footline}
{%
  \leavevmode%

  \begin{beamercolorbox}[wd=1\paperwidth,ht=2.5ex,dp=1.5ex,leftskip=.3cm plus1fill,rightskip=.3cm]{}%
    \usebeamerfont{title in head/foot}\insertframenumber/\inserttotalframenumber%
  \end{beamercolorbox}%
}

\setbeamertemplate{partpage}
{%
  \begin{centering}
    {\usebeamerfont{part name}\usebeamercolor[bg]{part name}\partname~\insertpartnumber}
    \begin{beamercolorbox}[sep=16pt,center]{part title}
      \usebeamerfont{part title}\insertpart\par
    \end{beamercolorbox}
  \end{centering}
}

\defbeamertemplate*{title page}{customised}[1][]
{%
    \usebeamerfont{title}\alert{\inserttitle}\par
    \usebeamerfont{subtitle}{\large\insertsubtitle}\par
    \bigskip
    \hrule
    \par
    \usebeamerfont{author}\insertauthor\par
    \usebeamerfont{institute}\small\insertinstitute\par
    \bigskip
    \hrule
    \par
    \usebeamerfont{date}{\normalsize\insertdate}\par
    \usebeamercolor[fg]{titlegraphic}\inserttitlegraphic{}
}

%%%

\setbeamertemplate{caption}[numbered]
\setbeamertemplate{caption label separator}{:}
\setbeamercolor{caption name}{fg=normal text.fg}

\usepackage{amssymb,amsmath}
\usepackage{mathtools}
\usepackage{fixltx2e} % provides \textsubscript

\usepackage{mathspec,xunicode,fontspec,xltxtra}
\defaultfontfeatures{Mapping=tex-text,Scale=MatchLowercase}
\setmainfont[BoldFont = Source Sans Pro Semibold, BoldItalicFont = Source Sans Pro Semibold Italic]{Source Sans Pro}
\setsansfont[BoldFont = Source Sans Pro Semibold, BoldItalicFont = Source Sans Pro Semibold Italic]{Source Sans Pro}
\setmonofont[Scale=0.8]{Source Code Pro}
\setmathfont(Digits,Latin,Greek)[Lowercase=Regular,Uppercase=Regular]{Clear Sans}

\usepackage{microtype}

% linguistics
\usepackage{tikz-qtree}
\usetikzlibrary{positioning,decorations.pathmorphing,decorations,trees,fit}
\tikzstyle{arrow} = [thin, <-, below left, bend right, font = \scriptsize]
\tikzset{>=latex}
\tikzset{every tree node/.style={align=center,anchor=north}}
\tikzset{execute at begin node=\vphantom{A}}
\newcommand*\circled[1]{\tikz[baseline=(char.base)]{
    \node[shape=circle, draw, inner sep=2pt] (char) {#1};}}

\usepackage{pgfplots}
\usepackage{tabulary}

\usepackage{expex}
\lingset{aboveglftskip=-1ex,
        belowglpreambleskip=-1ex,
        belowexskip=-1ex}

\resetcountonoverlays{excnt}

\usepackage{url}

\usepackage{graphicx}
\makeatletter
\def\maxwidth{\ifdim\Gin@nat@width>\linewidth\linewidth\else\Gin@nat@width\fi}
\def\maxheight{\ifdim\Gin@nat@height>\textheight0.8\textheight\else\Gin@nat@height\fi}
\makeatother
% Scale images if necessary, so that they will not overflow the page
% margins by default, and it is still possible to overwrite the defaults
% using explicit options in \includegraphics[width, height, ...]{}
\setkeys{Gin}{width=\maxwidth,height=\maxheight,keepaspectratio}

% Comment these out if you don't want a slide with just the
% part/section/subsection/subsubsection title:
\AtBeginPart{%
  \let\insertpartnumber\relax
  \let\partname\relax
  \frame{\partpage}
}
\AtBeginSection{%
  \let\insertsectionnumber\relax
  \let\sectionname\relax
  \frame{\sectionpage}
}
\AtBeginSubsection{%
  \let\insertsubsectionnumber\relax
  \let\subsectionname\relax
  \frame{\subsectionpage}
}

\usepackage[normalem]{ulem}
% avoid problems with \sout in headers with hyperref:
\pdfstringdefDisableCommands{\renewcommand{\sout}{}}

\setlength{\parindent}{0pt}
\setlength{\parskip}{6pt plus 2pt minus 1pt}
\setlength{\emergencystretch}{3em}  % prevent overfull lines

\setcounter{secnumdepth}{1}

\title{Agreement with possessed direct objects in Hungarian}

\subtitle{An experimental approach}

\author{\textsuperscript{a}András Bárány, \textsuperscript{b}Ádám Szalontai}

\date{SinFonIJA 8, Ljubljana, 25 September 2015}

\institute{\textsuperscript{a}ReCoS/DTAL, University of Cambridge,\\
        \textsuperscript{b}Research Institute for Linguistics,
        Hungarian Academy of Sciences}

\titlegraphic{%
    \begin{tikzpicture}[overlay]
        \node at (8.2, -2.5) {\includegraphics[height=2.5cm]{erc.eps}};
    \end{tikzpicture}
}

\begin{document}

\frame{%
\thispagestyle{empty} % no page number here
  \titlepage{}
}

\begin{frame}{Overview}
\tableofcontents[hideallsubsections]
\end{frame}

\section{Intro: Object agreement and possessed NPs}\label{case-and-case}

\begin{frame}{Object agreement in Hungarian}

\pex
    \a
        \begingl
            \gla    Mari lát egy kutyá-t.//
            \glb    M.      see.\textsc{3sg} a   dog-\textsc{acc}//
            \glft   ‘Mari sees a dog.’//
        \endgl
    \pause
    \a
        \begingl
            \gla    Mari lát-ja a kutyá-t.//
            \glb    M.   see-\textsc{\alert{3sg.obj}} the dog-\textsc{acc}//
            \glft   ‘Mari sees the dog.’//
        \endgl
\xe

    \begin{itemize}

        \item Roughly sensitive to definiteness of the DO

        \item Possessed direct objects are an exception

        \item Indefinite possessed DOs trigger agreement

        \item But: \alert{dialectal variation?}

    \end{itemize}

\end{frame}

\begin{frame}{This presentation}

    \begin{block}{Questions}

        \begin{itemize}

        \item What kind of variation do we find?

        \item What kinds of factors determine variation?

        \item What determines object agreement?

        \end{itemize}

    \end{block}

    \begin{block}{Methods and results}

        \begin{itemize}

            \item Series of online surveys

            \item Acceptability judgments with and without contexts

            \item Forced choice tests with and without contexts

        \end{itemize}
    \end{block}

\end{frame}

\section{The data}

\begin{frame}{Object agreement}

    \begin{itemize}

        \item Definiteness often a good predictor; not with possessed DOs

        \item Indefinite possessed DOs require agreement

    \end{itemize}

\onslide<3->{%
\pex
    \a
        \begingl
            \gla    Mari {lát /} \ljudge*lát-ja egy kutyá-t.//
            \glb    M.   see.\textsc{3sg} see-\textsc{\alert{3sg.obj}} a   dog-\textsc{acc}//
            \glft   ‘Mari sees a dog.’//
        \endgl
    \pause
    \a
        \begingl
            \gla    Mari {lát-ja /} \ljudge\%lát egy kutyá-m-at.//
            \glb    M.   see-\textsc{\alert{3sg.obj}} see.\textsc{3sg} the dog-\textsc{\alert{1sg.poss}-acc}//
            \glft   ‘Mari sees a dog of mine.’//
        \endgl
\xe}

\end{frame}

\begin{frame}{Types of possessed NPs}

\pex
    \a
        \begingl\rightcomment{Nominative possessor}
            \gla    Mari {(*a) /} egy kutyá-ja//
            \glb    Mari.\textsc{nom} (the) a dog-\textsc{3sg.poss}//
            \glft   `Mari's dog', `Mari's one dog'//
        \endgl
    \pause
    \a
        \begingl\rightcomment{Dative possessor}
            \gla    Mari-nak {a /} egy kutyá-ja//
            \glb    Mari-\textsc{dat} the a dog-\textsc{3sg.poss}//
            \glft   `Mari's dog', `a dog of Mari's'//
        \endgl
    \pause
    \a
        \begingl\rightcomment{Pronominal possessor}
            \gla    az ő kutyá-ja//
            \glb    the s/he.\textsc{nom} dog-\textsc{3sg.poss}//
            \glft   `her/his dog'//
        \endgl
\xe

\pause

\begin{itemize}

    \item Distribution suggests different heights in the NP
        \parencite{Szabolcsi1994,denDikken1999,Bartos1999,EKiss2000b,Dekany2015}

\end{itemize}

\end{frame}

\begin{frame}{Structures: \textsc{nom} possessor}

\ex\rightcomment{\emph{Mari kutyája} `Mari's dog'}
    \begin{tikzpicture}[baseline]
        \Tree   [.{\textcolor<2->{camcosblue}{DP}}
                    \node (specdp) {\textcolor<2->{camcosblue}{\emph{Mari.\textsc{nom}}}};
                    [.{\textcolor<2->{camcosblue}{D$'$}}
                        \textcolor<2->{camcosblue}{D}
                        [.PossP
                            \node (t) {\emph{\sout{Mari}}};
                                [.{Poss$'$}
                                    [.NP
                                        \node {\emph{kutyája}\\dog};
                                    ]
                                    [.Poss
                                        \node (poss) {\emph{-ja}\\\textsc{3sg.poss}};
                                    ]
                                ]
                        ]
                    ]
                ]
        \draw [arrow, ->, bend left=55] (t.south) to (specdp.south);
    \end{tikzpicture}
\xe

\end{frame}

\begin{frame}{Structures: \textsc{dat} possessor}

\ex\rightcomment{\emph{Marinak a kutyája} `Mari's dog'}
    \begin{tikzpicture}[baseline]
        \Tree   [.{\textcolor<2->{camcosblue}{DP}}
                    \node (specdp) {\textcolor<2->{camcosblue}{\emph{Marinak-\textsc{dat}}}};
                    [.{\textcolor<2->{camcosblue}{DP}}
                        \node (d) {\textcolor<2->{camcosblue}{D}\\\emph{a}\\the};
                        [.PossP
                            \node (t) {\emph{\sout{Mari}}};
                                [.{Poss$'$}
                                    [.NP
                                        \node {\emph{kutyája}\\dog};
                                    ]
                                    [.Poss
                                        \node (poss) {\emph{-ja}\\\textsc{3sg.poss}};
                                    ]
                                ]
                        ]
                    ]
                ]
        \draw [arrow, ->, bend left=55] (t.south) to (specdp.south);
    \end{tikzpicture}
\xe

\end{frame}

\begin{frame}{Structures: pronominal possessor}

\ex\rightcomment{\emph{az ő kutyája} `her/his dog'}
    \begin{tikzpicture}[baseline]
        \Tree   [.{\textcolor<2->{camcosblue}{DP}}
                    \node (d) {\textcolor<2->{camcosblue}{D}\\\emph{az}\\\emph{the}};
                    [.PossP
                        \node (t) {\textcolor<2->{camcosblue}{\emph{ő}}\\\textcolor<2->{camcosblue}{s/he}};
                            [.{Poss$'$}
                                [.NP
                                    \node {\emph{kutyája}\\dog};
                                ]
                                [.Poss
                                    \node (poss) {\emph{-ja}\\\textsc{3sg.poss}};
                                ]
                            ]
                    ]
                ]
    \end{tikzpicture}
\xe

\end{frame}

\begin{frame}{Syntactic and semantic properties}

    \begin{block}{Extraction}

        \begin{itemize}

            \item Dative possessors can be extracted \parencite{Szabolcsi1994}

            \item Only extracted possessors in definiteness effect-contexts (\nextx a)

            \item Non-specific reading in (\nextx a)

        \end{itemize}

    \end{block}

\onslide<2->{%
\pex
    \a
        \begingl
            \gla    Mari-nak van kutyá-ja.//
            \glb    Mari-\textsc{dat} is dog-\textsc{3sg.poss}//
            \glft   `Mari has a dog/dogs.'//
        \endgl
    \a
        \begingl
            \gla    \ljudge*Mari van kutyá-ja.//
            \glb    Mari.\textsc{nom} is dog-\textsc{3sg.poss}//
        \endgl
\xe}

\end{frame}

\begin{frame}{Variation}

    \begin{itemize}

        \item \textcite{Szabolcsi1994}: DOs with extracted possessors show
            variation

        \item \alert{Semantic agreement?} Only specific DOs agree

    \end{itemize}

\onslide<2->{%
\pex[*=\%]
    \a
        \begingl
            \gla    \ljudge\%Chomsky-nak nem olvas-t-ál vers-é-t.//
            \glb    Chomsky-\textsc{dat} \textsc{neg} read-\textsc{pst-2sg} poem-\textsc{3sg.poss-acc}//
            \glft   `You haven't read a poem of Chomsky's.'//
        \endgl
    \a
        \begingl
            \gla    Chomsky-nak nem olvas-t-ad vers-é-t.//
            \glb    Chomsky-\textsc{dat} \textsc{neg} read-\textsc{pst-\alert{2sg.obj}} poem-\textsc{3sg.poss-acc}//
            \glft   `You haven't read a poem of Chomsky's / Chomsky's poem.'//
        \endgl
\xe}

\end{frame}

\begin{frame}{What triggers object agreement?}

        \begin{itemize}

            \item Syntax? DP structure\hfill\parencite{Bartos1999}
            \item Semantics? [\textsc{def}] introduces presupposition\hfill\parencite{Coppock2013P}

        \end{itemize}

    \begin{block}{Goals of the surveys}

        \begin{itemize}

            \item Mapping the variation

            \item Finding out which factors determine variation \dots

            \item \dots{} and agreement

        \end{itemize}

    \end{block}

\end{frame}

\section{Surveys}

\begin{frame}{Methodology}

    \begin{itemize}

        \item Online surveys hosted on \texttt{http://spellout.net/ibexfarm/}

        \item April, May, September 2015: xxx participants

        \item Acceptability judgments (7 point Likert-scale):

            \begin{itemize}

                \item Subject vs.~object agreement

                \item Nominative vs.~dative possessors

                \item Positive vs.~negative polarity (semantic effect?)

                \item Word order (no effect?)

                \item Contexts: definite (unique) vs.~indefinite (non-unique) reading of the DO

            \end{itemize}

        \item Forced choice:

            \begin{itemize}

                \item Nominative vs.~dative possessors

                \item Presence vs.~absence of determiner

                \item Contexts: definite (unique) vs.~indefinite (non-unique) reading of the DO

            \end{itemize}

    \end{itemize}

\end{frame}

\begin{frame}{Data: Acceptability judgments}

    \pex In literature class, the children were reading Petőfi and Arany.
        \a
            \begingl
                \gla    Mari olvas-ott Petőfi vers-ét.//
                \glb    Mari read-\textsc{3sg.past} Petőfi.\textsc{nom} poem-\textsc{3sg.poss}//
            \endgl
        \a
            \begingl
                \gla    Mari olvas-t-a Petőfi vers-ét.//
                \glb    Mari read-\textsc{past-\alert{3sg.obj}} Petőfi.\textsc{nom} poem-\textsc{3sg.poss}//
            \endgl
        \a
            \begingl
                \gla    Mari olvas-ott Petőfi-nek vers-ét.//
                \glb    Mari read-\textsc{3sg.past} Petőfi-\textsc{\textcolor{red!70!black}{dat}} poem-\textsc{3sg.poss}//
            \endgl
        \a
            \begingl
                \gla    Mari olvas-t-a Petőfi-nek vers-ét.//
                \glb    Mari read-\textsc{past-\alert{3sg.obj}} Petőfi-\textsc{\textcolor{red!70!black}{dat}} poem-\textsc{3sg.poss}//
                \glft   \fbox{intended: `Mari read Petőfi's poem / a poem by Petőfi.'}//
            \endgl
    \xe

\end{frame}

\begin{frame}{Data: Forced choice I, no context}

%{\small
\hspace{-.5cm}\begin{tabulary}{1\linewidth}{ll}
Ismerte a szomszéd nővérét.        & Ismerte a szomszéd\textcolor{red!70!black}{nak} nővérét.\\
Ismerte a szomszéd nővérét.        & Ismerte a szomszéd\textcolor{red!70!black}{nak} \textcolor{green!70!black}{a} nővérét.\\
Ismerte a szomszéd nővérét.        & Ismerte a szomszéd\textcolor{red!70!black}{nak} \alert{egy} nővérét.\\
Ismerte a szomszéd nővérét.        & Ismerte a szomszéd \alert{egy} nővérét.\\
Ismerte a szomszéd \alert{egy} nővérét.    & Ismerte a szomszéd\textcolor{red!70!black}{nak} nővérét.\\
\multicolumn{1}{c}{\vdots} & \multicolumn{1}{c}{\vdots}\\
\end{tabulary}
%}

\begin{itemize}
    \item \emph{S/he knew the neighbour's daughter / \textcolor{green!70!black}{the} daughter \textcolor{red!70!black}{of the} neighbour / \alert{a} daughter \textcolor{red!70!black}{of the} neigbour / \dots}
\end{itemize}

\end{frame}

\begin{frame}{Data: Forced choice II, def.\ context}

\onslide<1->{%
\begin{itemize}
    \item \emph{\dots{} is a famous writer, but she also painted a landscape.}
\end{itemize}}

{\small
\hspace{-.5cm}\begin{tabulary}{1\linewidth}{ll}
Péter látta a költő \alert{egy} tájképét. & Péter látta a költő tájképet.\\
Péter látta a költő \alert{egy} tájképét. & Péter látta a költő\textcolor{red!70!black}{nek} tájképet.\\
Péter látta a költő \alert{egy} tájképét. & Péter látta a költő\textcolor{red!70!black}{nek} \alert{egy} tájképet.\\
Péter látta a költő \alert{egy} tájképét. & Péter látta a költő\textcolor{red!70!black}{nek} \textcolor{green!70!black}{a} tájképet.\\
Péter látta a költő tájképet.     & Péter látta a költő\textcolor{red!70!black}{nek} tájképet.\\
\multicolumn{1}{c}{\vdots} & \multicolumn{1}{c}{\vdots}\\
\end{tabulary}
}

\begin{itemize}
    \item \emph{Péter saw the painter \alert{one} landscape / the painter's landscape / \textcolor{green!70!black}{the} landscape \textcolor{red!70!black}{of the} painter / \dots}
\end{itemize}

\end{frame}

\begin{frame}{Data: Forced choice III, indef.\ context}

\begin{itemize}
    \item \emph{\dots{} is a famous writer, but she also painted \alert{several} landscapes.}
\end{itemize}

{\small
\hspace{-.5cm}\begin{tabulary}{1\linewidth}{ll}
Péter látta a költő \alert{egy} tájképét. & Péter látta a költő tájképet.\\
Péter látta a költő \alert{egy} tájképét. & Péter látta a költő\textcolor{red!70!black}{nek} tájképet.\\
Péter látta a költő \alert{egy} tájképét. & Péter látta a költő\textcolor{red!70!black}{nek} \alert{egy} tájképet.\\
Péter látta a költő \alert{egy} tájképét. & Péter látta a költő\textcolor{red!70!black}{nek} \textcolor{green!70!black}{a} tájképet.\\
Péter látta a költő tájképet.     & Péter látta a költő\textcolor{red!70!black}{nek} tájképet.\\
\multicolumn{1}{c}{\vdots} & \multicolumn{1}{c}{\vdots}\\
\end{tabulary}
}

\begin{itemize}
    \item \emph{Péter saw the painter \alert{one} landscape / the painter's landscape / \textcolor{green!70!black}{the} landscape \textcolor{red!70!black}{of the} painter / \dots}
\end{itemize}

\end{frame}

\section{Results}

\begin{frame}{Results}

    \begin{block}{Subject agreement with possessed DOs}

        \begin{itemize}

        \item Object agreement is always better than subject agreement
        \item Subject agreement has medians of 1
        \item blokk1\&2nomdat.png (exp 2), blokk1.png (exp 3)

        \item \alert{$\to$} We do not find judgments reported in the literature

        \end{itemize}

    \end{block}

    \begin{block}{What triggers agreement?}

        \begin{itemize}

        \item Not (in){}definiteness: possessed DOs can be indefinite, yet
            trigger agreement

        \item The possessive suffix does not influence definiteness

        \end{itemize}

    \end{block}

\end{frame}

\begin{frame}{Results II}

    \begin{itemize}

        \item \alert{$\to$} We do find significant effects between items with
            and without determiners

        \item Null possessors with and without determiners: \emph{néhány}
            `some', \emph{minden} `every' do not trigger agreement on their own

        \item B1/B2 (röviden)

        \item Dative possessors with and without determiners: \emph{egy} `a,
            one', \emph{a(z)} `the'

        \item DEF/IND/NIL, kontextus nélkül és kétféle kontextussal

    \end{itemize}

\end{frame}

\section{Conclusions}

\begin{frame}{Conclusions and future}

    \begin{itemize}

        \item We have not found data corroborating the variation reported

        \item All our speakers judged subject agreement with possessed objects
            as very low in acceptability

        \item The results highlight \alert{the role of determiners for acceptability}

        \item Why D\@?

        \begin{itemize}

            \item \textcite{Longobardi2008,Bernstein2008}: D as the syntactic
                locus of referentiality; DPs as arguments?

            \item Arguably explains extraction facts: no place for possessor?

        \end{itemize}

        \item Future: Non-specific possessed DOs? Other determiners?

    \end{itemize}

\end{frame}

\begin{frame}{Acknowledgements}

\begin{center}
    \Large \textbf{Thank you!}
\end{center}

The author's names are ordered by the size of their beards.  The first author's
research is supported by ERC Grant No. 269752 \enquote{Rethinking Comparative
    Syntax}. Thanks to: Theresa Biberauer, Anders Holmberg, Georg Höhn, Ian
Roberts, Michelle Sheehan, and Jenneke van der Wal.  The second author \dots

\end{frame}

\begin{frame}[t, allowframebreaks]

    \frametitle{References}
    \newrefcontext[sorting=nyt]
    \renewcommand*{\bibfont}{\footnotesize}
    \printbibliography[heading=none]

\end{frame}

\end{document}
